% ---------------------------------------------------------------------------
% Results.  Every table body below is written by make_results_figures.m.
% ---------------------------------------------------------------------------

All runs use the configuration of \cref{tab:config}.  The convergence study was
executed with \file{run_d1v5_gminmod_L2_convergence.m}.  The limiter comparison
and the $\theta$-sweep use $N_x = 6250$ with the same $\tau$, CFL and final time.

\subsection{Run summary and conservation}
\label{sec:conservation_results}

\Cref{tab:runs} confirms the time-step analysis of \cref{sec:dt}.  Every grid used
$\Delta t = \tau/4$, the same $120{,}001$ steps (the last is a round-off-sized remainder that lands
exactly on $t_{\mathrm{end}}$), and no retries or fallbacks.  \Cref{tab:conservation} checks the
conservation argument of \cref{sec:conservation}.  Mass and energy are conserved to
round-off relative to the discrete initial condition.  The momentum equals
$(p_L-p_R)\,t_{\mathrm{end}} = 0.135$ to round-off, because the undisturbed boundary
states exert the net pressure force $p_L - p_R$ on the domain.

\begin{table}[ht]
\centering
\caption{Convergence runs: grid spacing, time step, effective Courant number
$\nu_{\mathrm{eff}} = v_2\Delta t/\Delta x$, number of steps, and stability events.}
\label{tab:runs}
\small
\begin{tabular}{@{}rccccrr@{}}
\toprule
$N_x$ & $\Delta x$ & $\Delta t$ & $\nu_{\mathrm{eff}}$ & steps & retries & fallbacks \\
\midrule
\tabinput{data/tab_conv_setup.tex}
\bottomrule
\end{tabular}
\end{table}

\begin{table}[ht]
\centering
\caption{Conservation errors at $t=0.15$: $\sum_j q_j\Delta x$ minus its expected value
(mass and energy: discrete initial value; momentum: $(p_L-p_R)t$).}
\label{tab:conservation}
\small
\begin{tabular}{@{}rccc@{}}
\toprule
$N_x$ & mass & momentum & energy \\
\midrule
\tabinput{data/tab_conservation.tex}
\bottomrule
\end{tabular}
\end{table}

\paragraph{Odd grid.}  $N_x=3125$ is odd, so the cell centre $x_{1563}$ coincides with
the diaphragm.  The solver assigns cells with $x_j<x_0$ to the left state, so that
cell receives the right state.  The discrete diaphragm therefore sits at
$x_0-\Delta x/2$.  This is an $\ord(\Delta x)$ initial-condition offset present only
on the coarsest grid.  It contributes to that grid's contact position error
(\cref{tab:features}) and to the first pairwise orders.

\subsection{The computed solution}

\Cref{fig:finest} compares the $N_x = 50{,}000$ solution with the exact Euler
solution, and \cref{fig:gridzoom} shows how each wave sharpens under refinement.
The profiles are free of oscillations, and the four constant states are reproduced
to plotting accuracy.  The pointwise error is concentrated in a few cells at each
wave.  The largest errors occur at the shock in every variable and at the contact
in $\rho$ and $T$.  The rarefaction head and tail show small corner errors.  The
largest of these is a slight density undershoot just behind the tail
(\cref{fig:gridzoom}a), which is limited in amplitude and shrinks with $\Delta x$.
The maximum pointwise error does not decrease with $\Delta x$: it stays at
$\approx0.08$ in $\rho$ and $0.34$--$0.63$ in $u$.  This is expected, because an
$\ord(1)$ jump spread over a few cells always leaves an $\ord(1)$ error in the cell
that straddles it.

\begin{figure}[p]
\centering
\includegraphics[width=\textwidth]{fig_solution_finest.pdf}
\caption{Finest grid ($N_x=50{,}000$) versus the exact Euler solution at $t=0.15$
(left) and the pointwise error (right), inside the error window.}
\label{fig:finest}
\end{figure}

\begin{figure}[p]
\centering
\includegraphics[width=\textwidth]{fig_grid_zoom.pdf}
\caption{Grid refinement at the three waves.  Darker curves are finer grids.  The
shock (c) stays about three cells wide at every resolution.  The contact (b) is
wider in physical units and spans more cells as the grid is refined
(\cref{tab:features}).  The coarse-grid temperature overshoot in (d) disappears by
$N_x=12{,}500$.}
\label{fig:gridzoom}
\end{figure}

\subsection{Grid convergence}
\label{sec:conv_results}

\Cref{tab:L1,tab:L2,tab:orders} and \cref{fig:convergence} give the errors and
observed orders.  Because the solution outside the error window is exact to
round-off, the wave-region $L_2$ errors (\cref{tab:waveL2}) are the global ones
multiplied by $\approx\sqrt{L_x/0.7}=5.35$.  They therefore have identical orders.

\begin{table}[ht]
\centering
\caption{Global $L_1$ errors (mean absolute error over all cells).}
\label{tab:L1}
\small
\begin{tabular}{@{}rcccc@{}}
\toprule
$N_x$ & $\rho$ & $u$ & $p$ & $T$ \\
\midrule
\tabinput{data/tab_conv_L1.tex}
\bottomrule
\end{tabular}
\end{table}

\begin{table}[ht]
\centering
\caption{Global $L_2$ (RMS) errors and temperature peak excess.}
\label{tab:L2}
\small
\begin{tabular}{@{}rccccc@{}}
\toprule
$N_x$ & $\rho$ & $u$ & $p$ & $T$ & $T$ peak excess \\
\midrule
\tabinput{data/tab_conv_L2.tex}
\bottomrule
\end{tabular}
\end{table}

\begin{table}[ht]
\centering
\caption{Wave-region $L_2$ errors, $x\in[x_0-0.3,\ x_0+0.4]$.}
\label{tab:waveL2}
\small
\begin{tabular}{@{}rcccc@{}}
\toprule
$N_x$ & $\rho$ & $u$ & $p$ & $T$ \\
\midrule
\tabinput{data/tab_conv_waveL2.tex}
\bottomrule
\end{tabular}
\end{table}

\begin{table}[ht]
\centering
\caption{Observed orders.  Columns 2--5: pairwise $L_2$ orders~\eqref{eq:order} between
successive grids.  Then the least-squares $L_2$ order over all five grids, the
least-squares $L_1$ order, and the least-squares wave-region $L_2$ order.}
\label{tab:orders}
\small
\begin{tabular}{@{}lccccccc@{}}
\toprule
 & \multicolumn{4}{c}{pairwise $L_2$ order} & \multicolumn{3}{c}{least-squares fit} \\
\cmidrule(lr){2-5}\cmidrule(l){6-8}
 & $3125{\to}6250$ & $6250{\to}12.5$k & $12.5{\to}25$k & $25{\to}50$k & $L_2$ & $L_1$ & wave $L_2$ \\
\midrule
\tabinput{data/tab_conv_orders.tex}
\bottomrule
\end{tabular}
\end{table}

\begin{figure}[ht]
\centering
\includegraphics[width=\textwidth]{fig_convergence.pdf}
\caption{Convergence of the global $L_1$ and $L_2$ errors, with the $L_2$ error also
measured against the cell-averaged exact solution (right).  Dotted lines are
reference slopes.}
\label{fig:convergence}
\end{figure}

\paragraph{Interpretation.}  All four variables converge.  The least-squares $L_1$
orders are $0.81$--$0.89$, between the contact value $\tfrac23$ and the shock value
$1$ predicted in \cref{sec:expected_rates}.  For the variables that jump across the
contact ($\rho$, $T$) and for $u$, the $L_2$ orders are $0.40$--$0.49$, between the
contact value $\tfrac13$ and the shock value $\tfrac12$.  Pressure, which does not
jump at the contact, converges at $0.63$, somewhat above $\tfrac12$, because its
faster-converging rarefaction error still makes up a visible share of the total.  The ratio of the $L_2$ order to the $L_1$ order lies between $0.50$ and
$0.72$, close to the value of one half predicted by~\eqref{eq:L1L2delta}.  Temperature, which carries the largest relative contact jump, converges
slowest.  No pairwise order approaches two; the solution is not
smooth, and formal second-order accuracy cannot be verified on this test.

\subsection{Error by wave: which feature limits convergence}
\label{sec:features}

The global squared error is the sum of the contributions of the rarefaction, the
contact and the shock.  \Cref{fig:features} separates them using the zones
$[x_0-0.25,\,x_0+0.07]$, $[x_0+0.07,\,x_0+0.20]$ and $[x_0+0.20,\,x_0+0.40]$.
For density, the least-squares orders of the three contributions are
$\zoneSlopeAA$ (rarefaction), $\zoneSlopeAB$ (contact) and $\zoneSlopeAC$ (shock).
For temperature they are $\zoneSlopeDA$, $\zoneSlopeDB$ and $\zoneSlopeDC$.
The contact is the largest contributor on every grid and has the slowest
convergence, close to the theoretical $\tfrac13$.  It therefore controls the global
$\rho$ and $T$ errors, and its share grows under refinement.

\Cref{tab:features} and \cref{fig:widths} measure this directly.  The shock
thickness stays at $2.9$--$3.1$ cells (physical thickness order
$\shockWidthSlope\approx1$).  This is the self-steepening behaviour of a captured
shock.  The contact thickness scales as $\Delta x^{\contactWidthSlope}$, close to the
$\Delta x^{2/3}$ law for a second-order scheme, so it spans more cells on finer
grids: $4.3$ cells at $N_x=3125$ and $9.3$ cells at $N_x=50{,}000$.  The contact and
shock position errors decrease steadily.  The shock runs about $0.2$ cells ahead of
the exact position, and the contact lags by an amount that decreases faster than
$\Delta x$.

\begin{table}[ht]
\centering
\caption{Measured contact and shock thickness (maximum-gradient definition, from
$\rho$) in physical units and in cells, and the position errors
(numerical minus exact) of the mid-level crossings.}
\label{tab:features}
\small
\begin{tabular}{@{}rcccccc@{}}
\toprule
$N_x$ & $\delta_c$ & $\delta_c/\Delta x$ & $\delta_s$ & $\delta_s/\Delta x$ & contact position error & shock position error \\
\midrule
\tabinput{data/tab_features.tex}
\bottomrule
\end{tabular}
\end{table}

\begin{figure}[ht]
\centering
\includegraphics[width=\textwidth]{fig_feature_errors.pdf}
\caption{Contribution of each wave to the global $L_2$ error for (a) density and
(b) temperature.  The dashed band in (a) is the finite-$\tau$ contact floor
predicted by~\eqref{eq:erf_l2}.}
\label{fig:features}
\end{figure}

\begin{figure}[ht]
\centering
\includegraphics[width=0.62\textwidth]{fig_feature_widths.pdf}
\caption{Measured contact and shock thickness versus grid spacing.  The shaded band
is the kinetic contact thickness~\eqref{eq:dkin} predicted for
$\tau = 5\times10^{-6}$.}
\label{fig:widths}
\end{figure}

\subsection{The velocity error and the sub-cell position of the shock}
\label{sec:phase}

Velocity is the one variable whose $L_2$ error does not decrease monotonically:
$L_2(u)$ rises slightly from $N_x=3125$ to $6250$ and is flat from $25{,}000$ to
$50{,}000$ (\cref{tab:L2}).  This matches the ``velocity stops improving'' observation
recorded in earlier studies.  The zone decomposition shows why.  Between $93\%$
(coarsest grid) and $99\%$ (finest grid) of $L_2(u)^2$ comes from the $\pm4$ cells
around the shock, where $u$ jumps by $u^*=0.93$.  The velocity is continuous across
the contact, and the rarefaction contribution is four to ten times smaller than
the shock contribution.

The contribution from the handful of cells that straddle the shock depends on
\emph{where the exact shock falls inside a cell}.  The metric compares cell values
with the exact solution sampled at cell centres.  A shock near a face therefore
produces one cell with an almost full-jump error, while a shock near a cell centre
splits the error between two cells.  With the exact shock at
$x_s = x_0 + 0.262823$, its fractional position within its cell is
$0.57$, $0.13$, $0.26$, $0.53$ and $0.06$ on the five grids (\cref{tab:phase}).
Grouping the grids by this phase removes the irregularity.  Between the two
mid-cell grids ($N_x=3125$ and $25{,}000$) the shock-zone $L_2(u)$ converges with
order $\shockPhaseSlopeA$.  Between the two near-face grids ($6250$ and $50{,}000$)
the order is $\shockPhaseSlopeB$.  Both match the theoretical $\tfrac12$ for a captured
shock (\cref{fig:phase}a).  The flat $25{,}000\to50{,}000$ velocity result is
therefore a sampling artefact of comparing a point-sampled discontinuous reference
with an $\ord(1)$-per-cell error.  It does not indicate a relaxation-time or
time-step limit.

Averaging the exact solution over each cell, which is the natural reference for a
finite-volume method, reduces the coarse-grid velocity error.  It does not remove
the phase dependence (\cref{fig:phase}b, \cref{tab:orders_avg}), because the discrete
shock profile itself depends on the shock's sub-cell position.  Two practical
consequences follow.  First, $L_1(u)$, which weights the few shock cells less
heavily, converges monotonically with order $0.89$ and is the more robust velocity
metric.  Second, velocity $L_2$ orders should be quoted from least-squares fits over
several grids, preferably with refinement ratios that sample different shock phases,
rather than from a single pair of grids.

\begin{table}[ht]
\centering
\caption{Sub-cell position (phase) of the exact shock, the shock-zone ($\pm4\Delta x$)
velocity error, and that error scaled by $\Delta x^{1/2}$, against point-sampled and
cell-averaged exact solutions.  A phase-independent error would give a constant
scaled value.}
\label{tab:phase}
\small
\begin{tabular}{@{}rccccc@{}}
\toprule
 & & \multicolumn{2}{c}{point-sampled exact} & \multicolumn{2}{c}{cell-averaged exact} \\
\cmidrule(lr){3-4}\cmidrule(l){5-6}
$N_x$ & phase & shock-zone $L_2(u)$ & $/\Delta x^{1/2}$ & shock-zone $L_2(u)$ & $/\Delta x^{1/2}$ \\
\midrule
\tabinput{data/tab_shock_phase.tex}
\bottomrule
\end{tabular}
\end{table}

\begin{table}[ht]
\centering
\caption{Global $L_2$ errors against the cell-averaged exact solution, with pairwise
and least-squares orders (the last column is the least-squares $L_1$ order).}
\label{tab:orders_avg}
\small
\begin{tabular}{@{}rcccc@{}}
\toprule
$N_x$ & $\rho$ & $u$ & $p$ & $T$ \\
\midrule
\tabinput{data/tab_conv_L2avg.tex}
\bottomrule
\end{tabular}\\[6pt]
\begin{tabular}{@{}lcccccc@{}}
\toprule
 & $3125{\to}6250$ & $6250{\to}12.5$k & $12.5{\to}25$k & $25{\to}50$k & $L_2$ fit & $L_1$ fit \\
\midrule
\tabinput{data/tab_conv_orders_avg.tex}
\bottomrule
\end{tabular}
\end{table}

\begin{figure}[ht]
\centering
\includegraphics[width=\textwidth]{fig_shock_phase.pdf}
\caption{(a) Shock-zone velocity error.  Point labels give the sub-cell position of
the exact shock.  Grids with matching phase converge at order $\approx\tfrac12$.
(b) Global $L_2(u)$ against point-sampled and cell-averaged exact solutions.}
\label{fig:phase}
\end{figure}

\subsection{The role of the finite relaxation time}
\label{sec:tau_results}

\Cref{sec:contact_diffusion} predicts that at fixed $\tau$ the contact cannot become
thinner than $\delta_c^{\mathrm{kin}} = \kinContactMin$--$\kinContactMax$ for
$\tau=5\times10^{-6}$.  The contact error therefore has a floor of
$\rhoFloorMin$--$\rhoFloorMax$ in global $L_2(\rho)$.  The measured contact thickness
on the finest grid is $3.7\times10^{-3}$, about three times the kinetic value
(\cref{fig:widths}), and the contact-zone density error ($6.8\times10^{-4}$) is still
above the floor (\cref{fig:features}a).  The present study is therefore
discretization-dominated.  Extrapolating the measured $\Delta x^{0.73}$ contact
thickness, the numerical and kinetic thicknesses become equal near
$\Delta x\approx1\times10^{-4}$ ($N_x\approx2\times10^5$).  Below that spacing, further
refinement at $\tau=5\times10^{-6}$ would stop reducing the $\rho$ and $T$ errors.

For the earlier setting $\tau = 5\times10^{-5}$ the kinetic contact is
$10^{1/2}\approx3.2$ times thicker ($3.7$--$4.6\times10^{-3}$).  That is already equal
to the numerical contact thickness at $N_x = 50{,}000$.  This explains the plateau
recorded in the project handoff for that setting (\cref{tab:tau_compare}):
$L_2(\rho)$ was $1.130\times10^{-3}$ at $N_x=40{,}000$ and $1.132\times10^{-3}$ at
$50{,}000$, whereas at $\tau=5\times10^{-6}$ the error keeps falling to
$7.77\times10^{-4}$.  The earlier runs also used a ten times larger time step
($\Delta t = 1.25\times10^{-5}$, CFL $=0.10$), so the comparison does not separate
$\tau$ from $\Delta t$.  It is, however, quantitatively consistent with the
finite-$\tau$ contact floor.

\begin{table}[ht]
\centering
\caption{Global $L_2(\rho)$ at the two relaxation times.  The $\tau=5\times10^{-5}$ values
are the September 13 results recorded in the project handoff (CFL $=0.10$,
$\Delta t=1.25\times10^{-5}$).  The $N_x=50{,}000$ value at that $\tau$ comes from the
separate follow-on run in the same record.}
\label{tab:tau_compare}
\small
\begin{tabular}{@{}rccc@{}}
\toprule
$N_x$ & $\tau=5\times10^{-5}$ (handoff) & $\tau=5\times10^{-6}$ (this study) & ratio \\
\midrule
3,125  & $3.244\times10^{-3}$ & $3.163\times10^{-3}$ & $0.97$ \\
6,250  & $1.961\times10^{-3}$ & $1.806\times10^{-3}$ & $0.92$ \\
12,500 & $1.594\times10^{-3}$ & $1.392\times10^{-3}$ & $0.87$ \\
25,000 & $1.258\times10^{-3}$ & $1.003\times10^{-3}$ & $0.80$ \\
40,000 & $1.130\times10^{-3}$ & --- & --- \\
50,000 & $1.132\times10^{-3}$ & $7.766\times10^{-4}$ & $0.69$ \\
\bottomrule
\end{tabular}
\end{table}

\subsection{Limiter comparison}
\label{sec:limiter_results}

\Cref{tab:limiters} and \cref{fig:limiters} compare the five limiters developed
during the project at $N_x=6250$, under identical physics, $\Delta t$, grid and final
time.  The generalized-minmod case reproduces the convergence-study run at the same
grid to all printed digits, which cross-checks the two independently written
implementations.

\begin{table}[ht]
\centering
\caption{Limiter comparison, $N_x=6250$, $\tau=5\times10^{-6}$, CFL $=0.025$.}
\label{tab:limiters}
\footnotesize
\setlength{\tabcolsep}{4pt}
\begin{tabular}{@{}lcccccc@{}}
\toprule
Limiter & $L_1(\rho)$ & $L_1(u)$ & $L_1(p)$ & $L_1(T)$ & wave $L_1(T)$ & $T$ peak excess \\
\midrule
\tabinput{data/tab_limiters.tex}
\bottomrule
\end{tabular}
\end{table}

\begin{figure}[ht]
\centering
\includegraphics[width=\textwidth]{fig_limiters.pdf}
\caption{Limiter comparison at $N_x=6250$.  Minmod and the macroscopic hybrid nearly
coincide.  MC is the sharpest and shows the largest temperature overshoot, just
behind the contact (b).}
\label{fig:limiters}
\end{figure}

The results reproduce, at the new $\tau$ and CFL, the tradeoff identified earlier in
the project:
\begin{itemize}
  \item MC gives the smallest average errors but the largest temperature peak
        excess ($6.7\times10^{-3}$).  Its compressive slopes, $\phi$ up to $2$,
        overshoot at the contact, and the high-sensitivity populations of
        \cref{tab:sensitivity} amplify the overshoot into $T$.
  \item Minmod has the smallest overshoot ($4.5\times10^{-4}$) but the largest
        average errors.
  \item The macroscopic hybrid behaves almost exactly like minmod.  Its shared
        troubled-cell mask evidently covers most of the cells in which minmod and MC
        would differ, so it rarely uses the MC slope where it matters.
  \item Generalized minmod with $\theta=1.2$ reduces all four global $L_1$ errors by
        $20$--$23\%$ relative to minmod, and reduces MC's temperature peak excess by
        $77\%$.
\end{itemize}

\subsection{The $\theta$ sweep}
\label{sec:theta_results}

\Cref{tab:theta,fig:theta} show the whole family from $\theta=1$ (identical to
minmod) to $\theta=2$ (identical to MC).  The average error decreases monotonically
with $\theta$, and the temperature peak excess increases monotonically.  The family
therefore traces a single diffusion--overshoot tradeoff curve (\cref{fig:theta}b).
No $\theta$ is best on both metrics.  Measured along that curve, $\theta=1.2$ captures
$59\%$ of the $L_1(T)$ reduction available between minmod and MC while incurring
$18\%$ of the overshoot increase.  For comparison, $\theta=1.1$ gives $35\%/7\%$ and
$\theta=1.3$ gives $74\%/30\%$.  $\theta=1.2$ is thus a reasonable compromise, not an
optimum.  The appropriate value depends on how peak temperature error is weighed
against average error.  On the convergence grids the peak excess at $\theta=1.2$
vanishes for $N_x\ge12{,}500$ (\cref{tab:L2}), so the overshoot cost is a
coarse-grid effect at this $\tau$.  No fallback or retry occurred for any $\theta$.

\begin{table}[ht]
\centering
\caption{$\theta$ sweep, $N_x=6250$, $\tau=5\times10^{-6}$, CFL $=0.025$.  The last column
counts limiter fallbacks.}
\label{tab:theta}
\small
\begin{tabular}{@{}lccccccc@{}}
\toprule
$\theta$ & $L_1(\rho)$ & $L_1(u)$ & $L_1(p)$ & $L_1(T)$ & wave $L_1(T)$ & $T$ peak excess & fb \\
\midrule
\tabinput{data/tab_theta.tex}
\bottomrule
\end{tabular}
\end{table}

\begin{figure}[ht]
\centering
\includegraphics[width=\textwidth]{fig_theta_sweep.pdf}
\caption{(a) Global $L_1$ errors normalized by the minmod ($\theta=1$) values.
(b) Temperature $L_1$ error versus temperature peak excess.  Each point is one
$\theta$; the solver's $\theta=1.2$ is circled.}
\label{fig:theta}
\end{figure}
